\lysh{13368}{4}

\begin{alist}
\item Το εμβαδόν του $EZH\Theta$ είναι $(EZH\Theta) = (\Theta E)^2$. Στο ορθογώνιο τρίγωνο $AE\Theta$ είναι
$(AE) = x$ και $(A\Theta) = \alpha - x$. Από το πυθαγόρειο θεώρημα έχουμε:
$(\Theta E)^2 = (AE)^2 + (A\Theta)^2$
Άρα, $(EZH\Theta) = x^2 + (\alpha - x)^2$.
Επίσης, επειδή το $\Theta$ είναι σημείο της πλευράς $\Delta A$ και $x$ είναι η απόσταση από την κορυφή $\Delta$,
θα είναι:
\[0 \le (\Delta\Theta) \le (\Delta A) \Leftrightarrow 0 \le x \le \alpha.\]

\item Το εμβαδόν του $AB\Gamma\Delta$ είναι $(AB\Gamma\Delta) = \alpha^2$. Άρα η ζητούμενη σχέση γράφεται:
\begin{align*}
(EZH\Theta) \ge \dfrac{(AB\Gamma\Delta)}{2} \Leftrightarrow x^2 + (\alpha - x)^2 \ge \dfrac{\alpha^2}{2} \ &\Leftrightarrow \\
2x^2 + 2(\alpha - x)^2 \ge \alpha^2 \ &\Leftrightarrow \\
2x^2 + 2\alpha^2 - 4\alpha x + 2x^2 \ge \alpha^2 \ &\Leftrightarrow \\
4x^2 + \alpha^2 - 4\alpha x \ge 0 \Leftrightarrow (2x)^2 - 2\alpha(2x) + \alpha^2 \ge 0 \ &\Leftrightarrow
\end{align*}
$(2x - \alpha)^2 \ge 0$ που ισχύει.

\item Από το ερώτημα $\alpha)$ για $x = 1$ είναι:
\[(EZH\Theta) = 1^2 + (\alpha - 1)^2 = 1 + \alpha^2 - 2\alpha + 1 = \alpha^2 - 2\alpha + 2.\]
Οπότε η σχέση $(EZH\Theta) = \dfrac{2}{3}(AB\Gamma\Delta)$ ισοδύναμα γίνεται:
\begin{align*}
\alpha^2 - 2\alpha + 2 = \dfrac{2}{3}\alpha^2 \ &\Leftrightarrow \\
3\alpha^2 - 6\alpha + 6 = 2\alpha^2 \ &\Leftrightarrow \\
\alpha^2 - 6\alpha + 6 &= 0.
\end{align*}
Η εξίσωση είναι δευτέρου βαθμού ως προς $\alpha$ με διακρίνουσα $\Delta = (-6)^2 - 4 \cdot 6 = 12 > 0$ και
ρίζες:
\[\alpha_{1,2} = \dfrac{6 \pm \sqrt{12}}{2} = \dfrac{6 \pm 2\sqrt{3}}{2} = 3 \pm \sqrt{3} \Leftrightarrow \begin{cases} \alpha_1 = 3 - \sqrt{3} \approx 3 - 1,73 = 1,27 \\ \alpha_2 = 3 + \sqrt{3} \approx 4,73 \end{cases}\]
Από το ερώτημα $\alpha)$, πρέπει να ισχύει $x \le \alpha$, η οποία για $x = 1$ γίνεται $1 \le \alpha$.
Παρατηρούμε ότι και οι δύο τιμές είναι δεκτές.
\end{alist}

